\chapter{Processos, Troca de Contexto e Escalonamento}

\chapterauthor{João Pedro Rizzo Sales, Samuel Romano Ribeiro, Guilherme Michelsen do Monte}

\chapterrev{Ti Ribeiro Silvério}


\section{Descrição}


Este capítulo documenta três módulos do UFSKernel:
o módulo de \textbf{processos}, responsável por representar, criar e encerrar
unidades de execução; o mecanismo de \textbf{troca de contexto}, que realiza
a alternância segura entre modo usuário e modo kernel; e o \textbf{escalonador},
que decide qual processo recebe a CPU e por quanto tempo. Juntos, esses módulos
viabilizam a multiprogramação no kernel, conforme descrito em
Silberschatz et al.~\cite{silberschatz2018} e Tanenbaum e Bos~\cite{tanenbaum2015}.

Os arquivos que compõem esses módulos são: \texttt{process.h},
\texttt{process.c}, \texttt{proc\_ll.c}, \path{svc_entry.S}, \texttt{irq\_entry.S}, \texttt{irq\_return.S}, \texttt{context\_switch.S}, \texttt{scheduler.h} e \path{process_control_block.c}.


\section{Fundamentação Teórica}


\subsection{Processo e PCB}

Um processo é um programa em execução: uma entidade dinâmica que engloba
o contador de programa, os registradores, a pilha e o espaço de endereços
ocupado na memória~\cite{tanenbaum2015}. Para gerenciá-los, o kernel mantém
um \textbf{Bloco de Controle de Processo} (PCB) por processo, concentrando
tudo o que é necessário para suspendê-lo e retomá-lo como se nunca tivesse
sido interrompido~\cite{silberschatz2018}.

Durante seu ciclo de vida, um processo transita entre quatro estados. Os três
canônicos, \textbf{Pronto}, \textbf{Em execução} e \textbf{Bloqueado}, são
descritos em~\cite{tanenbaum2015}. O UFSKernel acrescenta um quarto,
\textbf{Zumbi}: o processo encerrou, mas sua estrutura ainda não foi
descartada, pois o processo pai deve registrar o encerramento do filho antes
que os recursos sejam liberados~\cite{silberschatz2018}.

\subsection{Trapframe, Contexto de Kernel e ABI ARM}

Na ARMv7, o processador opera em modos de privilégio distintos. Toda transição
de modo usuário para kernel, seja por chamada de sistema (\texttt{SVC}) ou por
interrupção (\texttt{IRQ}), exige salvar o estado completo do processo para que
a execução possa ser retomada. A estrutura que armazena
esse estado é o \textbf{trapframe}~\cite{tanenbaum2015}.

Para trocas de contexto internas ao kernel, quando o escalonador substitui
um processo por outro sem sair do modo supervisor, apenas um subconjunto
dos registradores precisa ser salvo. A ABI ARM (AAPCS)~\cite{arm:armv7ar} define
dois grupos: registradores \emph{call-preserved} (\texttt{r4}--\texttt{r11}, SP e LR), que uma
função deve restaurar antes de retornar, e registradores \emph{call-clobbered}
(\texttt{r0}--\texttt{r3}, \texttt{r12}), que o chamador sabe que podem ser destruídos. Como o compilador
já protege os call-clobbered em torno de qualquer chamada de função, o
mecanismo de troca de contexto de kernel precisa salvar apenas os call-preserved.

\subsection{Escalonamento: Round-Robin e MLFQ}

O escalonador decide qual processo pronto recebe a CPU e por quanto tempo
(o \textbf{quantum}), tipicamente acionado por um timer periódico~\cite{silberschatz2018}.
O \textbf{Round-Robin} organiza os processos em fila circular com quantum fixo,
mas trata igualmente processos interativos e CPU-bound~\cite{tanenbaum2015}.
A \textbf{Multilevel Feedback Queue} (MLFQ) supera essa limitação: mantém
múltiplas filas de prioridade com quanta crescentes, promovendo processos que
usam pouco CPU (interativos) e rebaixando os que esgotam o quantum (CPU-bound),
de modo que cada processo migra naturalmente para o nível compatível com seu
comportamento real~\cite{silberschatz2018}.


\section{Módulo de Processos}


\subsection{Decisões de Arquitetura e Design}

O módulo mantém duas estruturas distintas para o estado de execução de um
processo. O \texttt{struct trapframe} armazena todos os registradores de
propósito geral (\texttt{r0}--\texttt{r12}), os registradores banked do modo usuário
(\texttt{sp\_usr}, \texttt{lr\_usr}), o contador de programa (\texttt{pc\_usr})
e o registrador de estado (\texttt{cpsr\_usr}) — tudo o que o processo de
usuário enxerga. O \texttt{struct context} armazena apenas os registradores
call-preserved (\texttt{r4}--\texttt{r11}, SP e LR), suficientes para a troca de contexto
dentro do kernel. Salvar \texttt{r0}--\texttt{r3} e \texttt{r12} no contexto de kernel seria redundante,
pois o compilador já os trata como voláteis em torno de qualquer
chamada~\cite{arm:armv7ar}.

\subsubsection{Alocação Dinâmica de Memória}

As estruturas
fundamentais de um processo — o próprio PCB (\texttt{struct process}),
a pilha de usuário (\texttt{usr\_stack}) e a pilha de kernel (\texttt{kstack}) —
são alocadas dinamicamente requisitando blocos de 4~KB ao PMM através da 
função \texttt{pmm\_alloc\_block()}. O limite de 64 processos simultâneos existe para simplificar a gestão de PIDs, mas a memória física só é
efetivamente consumida sob demanda.

\subsubsection{Execução de Programas e Passagem de Argumentos}

Sem um sistema de arquivos externo disponível, os executáveis são representados
por funções compiladas juntamente com o kernel, gerenciadas pela
\texttt{program\_table} (definida em \path{programas.h}).

\lstset{captionpos=t}
\begin{lstlisting}[language=C, caption={Assinatura padronizada dos programas}]
void (*programa)(int argc, char** argv);
\end{lstlisting}
\lstset{captionpos=b}

Essa abordagem permite testar o ciclo completo \texttt{fork}--\texttt{exec}--\texttt{exit} 
simulando a passagem real de argumentos (\texttt{argc} e \texttt{argv}) para o 
espaço de execução do processo de forma análoga a sistemas POSIX reais.

\subsubsection{Gestão de PIDs e Bloqueio de Processos}

PIDs são gerenciados por \texttt{create\_pid} e \texttt{free\_pid} sobre
o array \path{used_pids[MAX_PROCESS_COUNT]}. O PID~0 é
reservado, convenção Unix que evita confundi-lo com um ponteiro
nulo~\cite{silberschatz2018}. O sistema emprega ativamente o estado de
bloqueio (\texttt{BLOCKED}). Um processo pode ser bloqueado ao aguardar por 
entrada de dados (E/S via UART) ou ao sincronizar com o término de um filho 
através da chamada \texttt{waitpid}, sendo reinserido na fila de prontos
assim que a condição de espera for satisfeita~\cite{tanenbaum2015}.

\subsection{Detalhes da Implementação}

\subsubsection{O PCB: \texttt{struct process}}

\lstset{captionpos=t}
\begin{lstlisting}[language=C, caption={Bloco de Controle de Processo em \texttt{process.h}}]
struct process {
    char* usr_stack;        // inicio da pilha de usuario
    char* kstack;           // inicio da pilha de kernel
    struct trapframe *tf;   // ponteiro para o trapframe (offset 8)
    enum State state;       // RUNNING, READY, BLOCKED ou ZOMBIE
    enum Block_Reason blocked_by;
    pid_t waiting_for_pid;  // PID do filho aguardado, se houver
    pid_t pid;
    process parent;
    pid_list children_ids;
    struct context context; // registradores call-preserved (troca de contexto)
    unsigned priority;
};
\end{lstlisting}
\lstset{captionpos=b}

A ordem dos três primeiros campos é estritamente deliberada: \texttt{usr\_stack} 
(4~bytes) e \texttt{kstack} (4~bytes) posicionam \texttt{tf} exatamente
no offset~8 — valor codificado na constante \texttt{TRAPFRAME\_OFFSET = 8},
amplamente utilizada pelo código assembly para acessar o trapframe do processo
diretamente.

\subsubsection{Ciclo de vida: \texttt{first\_process}, \texttt{fork}, \texttt{exec}, \texttt{exit} e \texttt{waitpid}}

\texttt{first\_process} cria o processo inicial: aloca as estruturas de pilha via PMM,
insere o processo no escalonador no nível de maior prioridade, define
\texttt{current} e invoca \texttt{sys\_exec} passando o ponteiro da função
solicitada.

\texttt{sys\_fork} cria um clone do processo. Ele requisita novas páginas ao
PMM para o \texttt{usr\_stack} e \texttt{kstack} do filho e copia o conteúdo
byte a byte da memória do pai. Os ponteiros do trapframe são meticulosamente
recalculados por aritmética de ponteiros baseada no deslocamento entre as pilhas:

\lstset{captionpos=t}
\begin{lstlisting}[language=C, caption={Ajuste de deslocamento na pilha no \texttt{sys\_fork}}]
unsigned int stack_diff = 
    (unsigned int)newprocess->usr_stack - (unsigned int)newprocess->parent->usr_stack;
newprocess->tf->sp_usr = newprocess->parent->tf->sp_usr + stack_diff;

unsigned int *saved_r7 = (unsigned int *)newprocess->tf->sp_usr;
*saved_r7 = *saved_r7 + stack_diff;
\end{lstlisting}
\lstset{captionpos=b}

O pai recebe o PID do filho em \texttt{\texttt{r0}}; o filho recebe zero, convenção
Unix clássica~\cite{silberschatz2018}. O \texttt{context.lr} do filho aponta
para \texttt{fork\_return\_asm}, garantindo que ele desembarque adequadamente
na restauração do modo usuário ao ser escalonado pela primeira vez.

\texttt{sys\_exec} prepara a imagem para a nova execução. Ele configura o 
\texttt{pc\_usr} para o ponteiro da nova função e repassa \texttt{argc} (em \texttt{\texttt{r0}}) e 
\texttt{argv} (em \texttt{r1}). Por motivos de isolamento e depuração, os 
registradores restantes (\texttt{r2} a \texttt{\texttt{r12}}) são zerados explicitamente. Em seguida, 
invoca-se \texttt{restore\_user\_context}.

\texttt{sys\_exit} marca o estado do processo atual como \texttt{ZOMBIE}. Crucialmente, 
verifica se o seu \texttt{parent} encontra-se bloqueado em \texttt{waitpid} aguardando
seu término (\texttt{waiting\_for\_pid}). Se afirmativo, muda o estado do pai para 
\texttt{READY}, efetivamente acordando-o. Por fim, invoca \texttt{pcb\_elect} 
para ceder imediatamente a CPU.

\texttt{sys\_waitpid} é o mecanismo responsável pela coleta (\emph{reaping}) de
processos Zumbis. Ao ser invocada, bloqueia o pai até que o processo filho 
especificado termine (mude para \texttt{ZOMBIE}). Uma vez encerrado, 
\texttt{waitpid} libera os blocos físicos de memória do filho de volta 
ao PMM, descarta seu PCB da tabela global, libera o PID na estrutura circular 
e retoma a execução do pai.

\section{Mecanismo de Troca de Contexto}


\subsection{Três mecanismos: entrada, retorno e troca}

O mecanismo de troca de contexto é composto por três partes complementares:
a \textbf{entrada} (\texttt{svc\_entry.S}, \texttt{irq\_entry.S}), que captura
o estado do processo de usuário e o armazena em um trapframe na pilha de kernel;
o \textbf{retorno} (\texttt{irq\_return.S}), que restaura esse estado e devolve
o controle ao processo; e a \textbf{troca de contexto de kernel}
(\texttt{context\_switch.S}), que salva e restaura apenas os registradores
call-preserved para alternar entre processos enquanto permanece em modo supervisor.
Mantê-los separados permite ao escalonador decidir, a cada interrupção, se
retorna ao mesmo processo ou realiza uma troca.

\subsection{Entrada via SVC e IRQ}

Ao receber um SVC (System Call) ou IRQ (Interrupção de Hardware), o processador ARMv7
salva automaticamente o CPSR em SPSR e o endereço de retorno em LR~\cite{arm:armv7ar}. 
A construção do trapframe, no entanto, deve ocorrer obrigatoriamente na pilha 
de kernel do processo (\texttt{SP\_svc}). 

Para chamadas de sistema (\texttt{svc\_entry.S}), o processador já ingressa 
no modo Supervisor, permitindo alocar o espaço e salvar o contexto de imediato.
Porém, para interrupções (\texttt{irq\_entry.S}), o processador entra no modo IRQ,
que possui sua própria pilha (\texttt{SP\_irq}). A solução empregada no UFSKernel é
salvar temporariamente os registradores mais voláteis na pilha IRQ, transitar 
manualmente para o modo SVC alterando o CPSR, alocar o trapframe na kstack e 
retornar brevemente ao modo IRQ para resgatar os valores originais:

\lstset{captionpos=t}
\begin{lstlisting}[language={[x86masm]Assembler},
                   caption={Transição de modo e resgate de registradores em \texttt{irq\_entry.S}}]
stmfd sp!, {\texttt{r0}-\texttt{r3}, lr}       @ Salva \texttt{r0}-\texttt{r3} e LR na pilha IRQ
mrs \texttt{r0}, cpsr
bic \texttt{r0}, \texttt{r0}, #0x1F
orr \texttt{r0}, \texttt{r0}, #0x13
msr cpsr_c, \texttt{r0}               @ Troca para o modo SVC (0x13)

sub sp, sp, #(17 * 4)        @ Aloca o trapframe na kstack (SP_svc)
add r1, sp, #(4 * 4)
stmia r1, {\texttt{r4}-\texttt{r12}}           @ Salva \texttt{r4}-\texttt{r12} em seguranca

@ ... (Transita de volta ao IRQ para ejetar a pilha e joga o LR em \texttt{r4})
\end{lstlisting}
\lstset{captionpos=b}

Após estabilizar a execução no modo Supervisor com os dados protegidos, 
finaliza-se a construção do trapframe:

\lstset{captionpos=t}
\begin{lstlisting}[language={[x86masm]Assembler},
                   caption={Finalização do trapframe na pilha de kernel}]
stmia sp, {\texttt{r0}-\texttt{r3}}            @ Salva \texttt{r0}-\texttt{r3} no inicio do trapframe
str \texttt{r4}, [sp, #(15 * 4)]      @ Salva o PC (resgatado em \texttt{r4})
str r1, [sp, #(16 * 4)]      @ Salva o SPSR original

add r2, sp, #(13 * 4)
stmia r2, {sp, lr}^          @ Salva SP_usr e LR_usr (^ acessa banked USR)
\end{lstlisting}
\lstset{captionpos=b}

O sufixo \texttt{\^} em \texttt{stmia r2, \{sp, lr\}\^} é essencial: ele
instrui o processador a acessar os registradores banked do modo usuário
(\texttt{SP\_usr} e \texttt{LR\_usr}), e não os do modo corrente. Sem ele,
o estado do processo seria corrompido~\cite{arm:armv7ar}. Antes de chamar
código C, a pilha é alinhada em 8 bytes, respeitando a convenção AAPCS.

\subsection{Retorno ao Modo Usuário}

\texttt{irq\_return.S} realiza o processo inverso: lê o trapframe no topo da
kstack, restaura o CPSR do usuário no SPSR, carrega \texttt{r0}--\texttt{r12}, restaura
\texttt{sp\_usr} e \texttt{lr\_usr} com o sufixo \texttt{\^}, ejeta o 
espaço da pilha e executa o retorno:

\lstset{captionpos=t}
\begin{lstlisting}[language={[x86masm]Assembler},
                   caption={Retorno ao modo usuário em \texttt{irq\_return.S}}]
ldr \texttt{r0}, [sp, #(16 * 4)]      @ \texttt{r0} = SPSR salvo
msr spsr_cxsf, \texttt{r0}            @ joga para SPSR_svc

ldr lr, [sp, #(15 * 4)]      @ carrega o PC salvo para LR_svc
add \texttt{r0}, sp, #(13 * 4)
ldmia \texttt{r0}, {sp, lr}^          @ restaura SP_usr e LR_usr
nop                              
ldmia sp, {\texttt{r0}-\texttt{r12}}           @ restaura \texttt{r0}-\texttt{r12}
add sp, sp, #(17 * 4)        @ Ejetamos o trapframe da kstack!

movs pc, lr                  @ Pula para o PC e restaura o CPSR nativamente
\end{lstlisting}
\lstset{captionpos=b}

O sufixo \texttt{s} na instrução \texttt{movs} com o PC como destino em 
modo privilegiado copia automaticamente o SPSR para o CPSR~\cite{arm:armv7ar} — 
restaurando modo, flags de processador e conjunto de instruções em uma única 
operação atômica.

\subsection{Troca de Contexto de Kernel}

\lstset{captionpos=t}
\begin{lstlisting}[language={[x86masm]Assembler},
                   caption={Troca de contexto de kernel em \texttt{context\_switch.S}}]
context_switch:
    cmp \texttt{r0}, #0
    beq if_null              @ se nao ha processo saindo, apenas carrega

    stmia \texttt{r0}!, {\texttt{r4}-\texttt{r11}}      @ salva registradores call-preserved
    str sp, [\texttt{r0}], #4         @ salva SP
    str lr, [\texttt{r0}]             @ salva LR (ponto de retorno quando reativado)

if_null:
    ldmia r1!, {\texttt{r4}-\texttt{r11}}      @ restaura registradores call-preserved
    ldr sp, [r1], #4         @ restaura SP (agora na pilha do novo processo)
    ldr lr, [r1]             @ restaura LR

    bx lr                    @ retorna -- no contexto do processo entrante
\end{lstlisting}
\lstset{captionpos=b}

O caso \texttt{\texttt{r0} == NULL} cobre o boot inicial, quando não há um 
processo anterior saindo da CPU. A propriedade central do \texttt{context\_switch} 
é que o \texttt{bx lr} não retorna para quem chamou a função de troca, mas para 
o ponto exato em que o processo entrante havia sido suspenso. Para um processo 
recém-criado via \texttt{fork}, esse ponto de entrada é forjado como 
\texttt{fork\_return\_asm} (que redireciona para \texttt{irq\_return}), 
fazendo com que o novo fluxo de execução desembarque diretamente e de 
forma segura no modo usuário.

\section{Escalonador}

\subsection{Decisões de Arquitetura e Design}

\subsubsection{MLFQ com 32 níveis e quanta crescentes}

A ideia adotada é uma MLFQ simplificada com 32 níveis. O \texttt{pcb}
é um array de \texttt{struct pcb\_entry} com quanta variando de 20 a 175
em progressão aritmética de passo 5. Processos recém-criados entram no
nível~0, maior prioridade, menor quantum, refletindo a heurística de que
processos novos tendem a ser interativos~\cite{tanenbaum2015}. A progressão
aritmética produz separação gradual entre os níveis sem saltos abruptos que
pudessem causar \emph{starvation} nos níveis inferiores.

\subsubsection{Margem de quantum e \texttt{list\_location}}

A constante \texttt{SCHEDULER\_QUANTUM\_MARGIN = 5} define uma zona de
tolerância: sem ela, pequenas variações no tempo acumulado causariam
oscilações de promoção e rebaixamento. A variável global \texttt{list\_location}, atualizada por
\texttt{pcb\_get\_node\_pid}, comunica o nível atual
do processo para \texttt{pcb\_climb} e \texttt{pcb\_fall}. É uma escolha que evita parâmetros extras, mas cria acoplamento implícito:
uma chamada intercalada a \texttt{pcb\_get\_node\_pid} com outro PID
sobrescreveria \texttt{list\_location} e corromperia a movimentação. Em um
sistema multiprocessado, isso seria uma fonte clássica de \emph{race conditions}.

\subsubsection{Alocação estática e \texttt{time\_elapsed}}

Os nós da lista ligada (que organizam as filas do PCB) são alocados estaticamente em
\texttt{process\_blocks\_nodes[64]}, indexados por PID, impondo o mesmo
limite de 64 processos. O comportamento do escalonador é
funcional, mas vale ressaltar que a função \texttt{time\_elapsed} atualmente simula a 
passagem de tempo, e o feedback de prioridade não reflete ciclos precisos de relógio real.

\subsection{Detalhes da Implementação}

\subsubsection{Estruturas de dados}

\lstset{captionpos=t}
\begin{lstlisting}[language=C, caption={Estruturas do escalonador em \texttt{scheduler.h}}]
struct pcb_entry {
    pll_node *process_list; // cabeca da lista ligada de processos
    unsigned quantum;       // quantum deste nivel
    short process_count;    // numero de processos no nivel
    short next_process;     // indice circular do proximo a ser eleito
};

typedef struct process_node_ll {
    struct process_node_ll *next;
    process proc;
} pll_node;
\end{lstlisting}
\lstset{captionpos=b}

O campo \texttt{next\_process} implementa o revezamento circular dentro de
cada nível: após eleger o processo de índice $i$, o campo avança para
$(i+1) \bmod \texttt{process\_count}$.

\subsubsection{O algoritmo de eleição: \texttt{pcb\_elect}}

Invocada a cada \emph{tick} do timer via \texttt{timer\_callback}, ou sempre que 
um processo cede a CPU voluntariamente (ex: chamando \texttt{waitpid} ou \texttt{exit}),
a função \texttt{pcb\_elect} executa em três fases:

\textbf{Fase 1 — Acumular tempo.} Adiciona o retorno de \texttt{time\_elapsed()}
à variável \path{running_for}.

\textbf{Fase 2 — Avaliar o processo atual.} Se o processo ativo ainda está dentro
do quantum (com a margem) e não mudou de estado, a função retorna sem preemptar:

\lstset{captionpos=t}
\begin{lstlisting}[language=C, caption={Avaliação do quantum em \texttt{pcb\_elect}}]
if (current->state == RUNNING) {
    if (running_for < (quantum + SCHEDULER_QUANTUM_MARGIN))
        return; // ainda dentro do quantum, mantem rodando

    // esgotou o quantum: aplica feedback de prioridade
    if (running_for <= quantum - SCHEDULER_QUANTUM_MARGIN)
        pcb_climb(current_process_node); // usou pouco: promove
    else
        pcb_fall(current_process_node);  // usou muito: rebaixa
        
    running_for = 0;
    current->state = READY;
    current->blocked_by = BT_TIMER;
}
\end{lstlisting}
\lstset{captionpos=b}

\textbf{Fase 3 — Eleger o próximo.} Percorre os 32 níveis do PCB do mais alto para
o mais baixo (nível 0 ao 31). Em cada nível que contenha processos (\texttt{process\_count > 0}),
um laço interno avalia as opções disponíveis garantindo que não entre em loop infinito caso
haja inconsistências:

\lstset{captionpos=t}
\begin{lstlisting}[language=C, caption={Laço de tentativa e eleição de processo}]
for(int tries = 0; tries < pcb[i].process_count; tries++) {
    pll_node *current_node = pll_idx(pcb[i].process_list, pcb[i].next_process);

    if(current_node != NULL && current_node->proc->state == READY) {
        process old_current = current;
        running_for = 0;
        pcb[i].next_process = (pcb[i].next_process + 1) % pcb[i].process_count;

        current = current_node->proc;
        current->state = RUNNING;
        context_switch(&old_current->context, &current->context);
        return;
    }
    pcb[i].next_process = (pcb[i].next_process + 1) % pcb[i].process_count;
}
\end{lstlisting}
\lstset{captionpos=b}

Este laço de tentativas (\texttt{tries}) checa circularmente a fila. 
Ao encontrar um candidato válido em estado \texttt{READY}, atualiza os 
índices de \texttt{current}, altera o estado do novo eleito, 
invoca \texttt{context\_switch} e encerra imediatamente a eleição.

\subsubsection{Promoção e rebaixamento}

\texttt{pcb\_climb} e \texttt{pcb\_fall} movem o processo um nível acima
ou abaixo, respectivamente, removendo o nó do nível atual e inserindo-o no
adjacente. Os limites (nível~0 e nível~31) são respeitados rigidamente: processos nos
extremos não são movidos além deles.


\section{Interface Pública (API)}


\subsection{Módulo de Processos (\texttt{process.h})}

\begin{description}
    \item[\texttt{pid\_t create\_pid()}] Aloca e retorna um PID disponível.
    Retorna \texttt{MAX\_PROCESS\_COUNT} se todos estiverem em uso.

    \item[\texttt{void free\_pid(pid\_t pid)}] Libera o PID para reutilização.

    \item[\texttt{void first\_process(void (*programa)(int, char**))}] Cria e lança
    o processo inicial, passando o ponteiro da função a ser executada como ponto
    de entrada. Deve ser chamada uma única vez pelo \texttt{kmain}. Não retorna.

    \item[\texttt{int fork()}] Cria um processo filho. Retorna o PID do filho
    ao pai e 0 ao filho (wrapper assembly em \texttt{syscall.S}).

    \item[\texttt{void exec(void (*programa)(int, char**), int argc, char** argv)}] Substitui a imagem do
    processo atual pelo novo programa repassando seus argumentos (\texttt{argc} e \texttt{argv}). Não retorna ao chamador.

    \item[\texttt{int waitpid(int pid)}] Bloqueia o processo chamador (\texttt{BLOCKED}) até que o processo filho 
    especificado pelo PID termine, garantindo a coleta do estado Zumbi e liberação da memória. Retorna o PID finalizado.

    \item[\texttt{void exit()}] Encerra o processo atual e avisa ao pai, se este estiver esperando. Não retorna.
\end{description}

\subsection{Módulo do Escalonador (\texttt{scheduler.h})}

\begin{description}
    \item[\texttt{void pcb\_elect(void)}] Seleciona e ativa o próximo
    processo. Chamada automaticamente pelo timer; pode ser invocada
    explicitamente por bloqueios de I/O, \texttt{sys\_waitpid} ou \texttt{sys\_exit}.

    \item[\texttt{void pcb\_insert(unsigned priority\_level, pll\_node* proc)}]
    Insere um processo no nível especificado. Nível~0 é o de maior prioridade.

    \item[\texttt{void pcb\_remove(unsigned priority\_level, pll\_node* proc)}]
    Remove um processo do nível especificado.

    \item[\texttt{pll\_node* pcb\_get\_node\_pid(pid\_t pid)}] Localiza o nó
    do processo pelo PID. Efeito colateral: atualiza \texttt{list\_location}.

    \item[\texttt{void pcb\_climb(pll\_node* proc)}] Move o processo um nível
    acima. Requer chamada prévia de \texttt{pcb\_get\_node\_pid} para o
    mesmo processo.

    \item[\texttt{void pcb\_fall(pll\_node* proc)}] Move o processo um nível
    abaixo. Mesma precondição de \texttt{pcb\_climb}.

    \item[\texttt{pll\_node* pll\_node\_new(process proc)}] Inicializa e
    retorna o nó de lista ligada alocado estaticamente para o processo,
    indexado pelo PID.

    \item[\texttt{int proc\_is\_valid(process proc)}] Retorna verdadeiro se
    o processo não é nulo e não está em estado \texttt{ZOMBIE}.
\end{description}

%-------------------------------------------------------------------------------
\section{Testes e Verificação}
%-------------------------------------------------------------------------------

\subsection{Rotina de Teste}

O \texttt{kmain} inicializa os gerenciadores de memória 
(PMM e VMM) e invoca \path{first_process(prog_init_idle)}. Este processo 
primário cria um filho via \texttt{fork} e executa (\texttt{exec}) a 
\texttt{prog\_shell}, cedendo controle ao usuário através de uma interface interativa.

O teste de validação do módulo de processos e escalonamento consiste em interagir
com essa shell e invocar o comando \texttt{multitask}. Esse programa invoca 
sucessivos \texttt{forks} para criar três processos filhos simultâneos que executam
laços de espera ocupada. Enquanto isso, o processo pai 
utiliza \texttt{waitpid} para aguardar o término sequencial de cada filho. 

\subsection{Saída Esperada}

Ao executar \texttt{make run} e interagir com o sistema, a saída no console 
será análoga à listagem abaixo:

\lstset{captionpos=t}
\begin{lstlisting}[caption={Saída serial do QEMU executando \texttt{multitask}}]
> multitask
Processo pai iniciando...
Filho 1: PID 4
Filho 2: PID 5
Filho 3: PID 6
Pai esperando os filhos...
Esperando 1: PID 4
Filho 1: PID 0
   -> Filho 1 nasceu!
Filho 3: PID 0
   -> Filho 3 nasceu!
Filho 2: PID 0
   -> Filho 2 nasceu!
   <- Filho 1 finalizou!
Esperando 2: PID 5
   <- Filho 3 finalizou!
   <- Filho 2 finalizou!
Esperando 3: PID 6
Todos os filhos terminaram. Pai saindo.
> exit
\end{lstlisting}
\lstset{captionpos=b}

\subsection{Análise da Execução}

A saída comprova o funcionamento integrado de múltiplos subsistemas do kernel:

\begin{enumerate}
    \item \textbf{Criação de Processos (\texttt{fork} e \texttt{exec}):} A shell
    consegue interpretar o comando, realizar o \texttt{fork} e substituir a imagem
    do processo filho pelo código do \texttt{prog\_multitask}. O processo pai, por sua 
    vez, cria três filhos (PIDs 4, 5 e 6).
    
    \item \textbf{Escalonamento e Preempção:} A mensagem \texttt{"Esperando 1: PID 4"} 
    é imediatamente seguida pelas inicializações dos três filhos intercaladas 
    (\texttt{"Filho 1 nasceu!"}, \texttt{"Filho 3 nasceu!"}, \texttt{"Filho 2 nasceu!"}). 
    Isso prova que o timer está interrompendo a CPU ativa e o escalonador MLFQ 
    está realizando as trocas de contexto com sucesso entre os três processos irmãos.
    
    \item \textbf{Sincronização e Coleta (\texttt{waitpid}):} O pai bloqueia em 
    \texttt{waitpid(4)}. Quando o filho 1 finaliza (\texttt{"<- Filho 1 finalizou!"}), 
    o pai é acordado e imediatamente emite \texttt{"Esperando 2: PID 5"}. O processo se repete
    até que todos os filhos terminem e o pai encerre, retornando o controle limpo
    à shell (prompt \texttt{>}), comprovando que a memória e os PIDs foram
    integralmente desalocados sem vazamentos.
\end{enumerate}
